\documentclass[11pt,openany]{book}

\usepackage[
    left=2.7cm,
    right=2.7cm,
    top=3cm,
    bottom=3cm,
    headheight=14pt
]{geometry}

\usepackage[T1]{fontenc}
\usepackage{lmodern}
\usepackage{microtype}

\usepackage{amsmath}
\usepackage{amssymb}
\usepackage{amsthm}
\usepackage{mathtools}

\usepackage{tikz}
\usepackage{float}
\usepackage{graphicx}
\usepackage{booktabs}
\usepackage{xcolor}
\usepackage{enumitem}
\usepackage{caption}
\usepackage{titlesec}
\usepackage{fancyhdr}
\usepackage[most]{tcolorbox}
\usepackage{hyperref}

\hypersetup{
    colorlinks=true,
    linkcolor=blue!50!black,
    urlcolor=blue!50!black,
    citecolor=blue!50!black,
    pdftitle={Road to Relativity},
    pdfauthor={Nishant Mathew}
}

\setlength{\parindent}{0pt}
\setlength{\parskip}{0.65em}
\setlist{itemsep=0.25em, topsep=0.35em}

\pagestyle{fancy}
\fancyhf{}
\fancyhead[LE,RO]{\thepage}
\fancyhead[LO]{\nouppercase{\rightmark}}
\fancyhead[RE]{\nouppercase{\leftmark}}
\renewcommand{\headrulewidth}{0.4pt}

\titleformat{\chapter}[display]
    {\normalfont\huge\bfseries}
    {\chaptername\ \thechapter}
    {20pt}
    {\Huge}

\titleformat{\section}
    {\normalfont\Large\bfseries}
    {\thesection}
    {0.9em}
    {}

\titleformat{\subsection}
    {\normalfont\large\bfseries}
    {\thesubsection}
    {0.8em}
    {}

\theoremstyle{definition}
\newtheorem{definition}{Definition}[chapter]
\newtheorem{example}{Example}[chapter]
\newtheorem{remark}{Remark}[chapter]

\tcbset{
    enhanced,
    breakable,
    arc=2mm,
    boxrule=0.5pt,
    left=7pt,
    right=7pt,
    top=7pt,
    bottom=7pt
}

\newtcolorbox{keyidea}{
    title=\textbf{Key Idea},
    colback=blue!5,
    colframe=blue!50!black
}

\newtcolorbox{definitionbox}{
    title=\textbf{Definition},
    colback=green!5,
    colframe=green!40!black
}

\newtcolorbox{warningbox}{
    title=\textbf{Common Mistake},
    colback=red!5,
    colframe=red!50!black
}

\newtcolorbox{examplebox}{
    title=\textbf{Example},
    colback=yellow!8,
    colframe=yellow!40!black
}

\title{\Huge Mathematics for the Physically Inclined\\[0.5em]
\Large A Curious Journey into the Mathematics of Advanced Physics}
\author{Nishant Mathew}
\date{April 2026}

\begin{document}

\frontmatter

\maketitle


\vspace*{0.3\textheight}

\begin{quote}
\centering
``Curiouser and curiouser!''

\medskip

--- Lewis Carroll, \textit{Alice's Adventures in Wonderland}
\end{quote}

\vspace{1cm}





\chapter*{Preface}
\addcontentsline{toc}{chapter}{Preface}

Mathematics and physics are often introduced as collections of formulas, procedures, 
and abstract definitions to be memorised. Students learn to manipulate symbols long 
before they are shown what those symbols mean geometrically or physically. As a result, 
ideas that are often simple at their core can begin to feel mysterious, disconnected, 
or purely mechanical.

\medskip

This book approaches these subjects differently.


\medskip

The goal of this text is not merely to teach linear algebra, tensors, or relativity as 
computational subjects, but to develop an intuitive understanding of the structures that 
underlie them. Throughout the book we will repeatedly return to simple geometric ideas 
--- walking on a hill, changing coordinate systems, stretching space, comparing directions 
--- and use them to uncover the mathematics beneath.

\medskip

The philosophy of this book is that mathematics should arise naturally from geometry and 
physical intuition. Rather than beginning with formal definitions and postponing meaning 
until later, we will often proceed in the opposite direction: beginning with geometric 
ideas first, and then carefully building the formal mathematics needed to describe them.

\medskip

A recurring theme throughout the text is the distinction between an object and its 
description. A vector is not its coordinates. A transformation is not the matrix used 
to represent it. A geometric or physical statement should not depend on the coordinate 
system chosen to describe it. This idea leads naturally toward vectors, covectors, 
metrics, tensors, and eventually the geometry of spacetime itself.

\medskip

Like Alice following the rabbit down the rabbit hole, the reader is invited to begin with something familiar and simple, only to discover that it leads into a much deeper world.
from a desire to understand these subjects conceptually rather than purely computationally. 
Many of the explanations in these pages arose from repeatedly asking a simple question:

\begin{quote}
What do these mathematical objects actually \textit{mean} geometrically?
\end{quote}

\medskip



This book is therefore an invitation to move slowly, think geometrically, and follow the 
rabbit hole patiently --- allowing familiar ideas to become strange, and then gradually 
become clear again.
\tableofcontents

\mainmatter

\part{Geometry and Motion}



\vspace*{0.35\textheight}

\begin{quote}
\centering
``Would you tell me, please, which way I ought to go from here?''

\medskip

``That depends a good deal on where you want to get to,'' said the Cat.

\medskip

--- Lewis Carroll, \textit{Alice's Adventures in Wonderland}
\end{quote}


Alice's journey begins with a question of direction. In Wonderland, paths are strange, 
size is uncertain, and familiar things no longer behave in familiar ways. To move through 
such a world, Alice must first learn how to understand where she is, where she is going, 
and what it means to change perspective.

\medskip

This part of the book begins in the same spirit. We start with motion from one place to 
another, and from this simple idea build the language of vectors, coordinates, bases, 
linear transformations, and change of basis.

\medskip

The central theme is the distinction between an object and its description. A vector is 
not its coordinates. A transformation is not merely its matrix. When we change our basis, 
the numbers may change, but the geometric object remains the same.

\medskip

This is the first step down the rabbit hole: from ordinary motion to the geometry beneath 
physics.


\chapter{Vectors and Coordinates}

\section{What Is a Vector Really?}

Suppose you are running up a hill. You start at the bottom, at Point A, and end somewhere in the middle of the hill, at Point B. We can represent the overall change in position with a diagram.

\begin{figure}[htbp]
\centering
\begin{tikzpicture}
  \draw[thick] (-6,0) -- (0,4) -- (0,0) -- (-6,0);
  \draw[->, red, very thick] (-6,0) -- (-3,2) node[right] {$\vec{v}$};
  \fill (-6,0) circle (1.5pt) node[left]{A};
  \fill (-3,2) circle (1.5pt) node[above=2pt]{B};
\end{tikzpicture}
\caption{Vector $\vec{v}$ from point A to B.}
\label{fig:vector_diagram}
\end{figure}

\textit{It is important to note that the arrow does not represent the exact path taken by the runner. The runner may have taken a curved or uneven route. }

\medskip

The vector only captures the overall change from A to B, ignoring the details of the journey.

\medskip

A vector is often described as an object that has magnitude and direction. That is, it tells us how far something goes, and in what direction.

\medskip

Now consider the red arrow in the diagram. It points from A to B, and it clearly has a certain length. 
So it seems to fit this description. But there is something subtle going on. The diagram does not tell us what we mean by directions such as left, right, up, or down in any precise way. 
It also does not tell us how to measure length. We have not said what counts as one unit of distance.

\medskip

Does this mean the red object is not a vector? No.

\medskip

The vector is still there. It represents the displacement from A to B. It tells us how to move from the starting point to the ending point. What is missing is not the vector itself, but a way of describing it using numbers.

\medskip

If we decide how to measure distances, and if we agree on how to describe directions, 
then we can assign numbers to this vector. We could then say exactly how long it is, 
and describe its direction more precisely.

\medskip

So the vector itself does not depend on these choices, but any numerical description of it does.

\medskip

This raises an important question: how do we actually describe a vector using numbers?

\medskip

To do that, we will need to introduce a way of measuring lengths and describing directions more precisely.

\section{Describing Direction and Length}

In the previous section, we saw that a vector represents a displacement from one point to another. 
It tells us how to move, but we have not yet said how to describe that movement using numbers.

\medskip

If we want to be more precise, we need two things:

\begin{enumerate}
  \item a way to measure distance,
  \item a way to describe direction.
\end{enumerate}

\medskip

Let us start with distance. We shall introduce a grid to our previous diagram.

\begin{figure}[H]
\centering
\begin{tikzpicture}[scale=0.8]

  \draw[step=1cm, gray!30, thin] (-7,-1) grid (2,5);

  \draw[->, thick] (-7,0) -- (2.5,0) node[right] {$x$};
  \draw[->, thick] (0,-1) -- (0,5.5) node[above] {$y$};

  \foreach \x in {-6,-5,-4,-3,-2,-1,1,2}
    \draw (\x,0.1) -- (\x,-0.1) node[below] {\small $\x$};

  \foreach \y in {1,2,3,4,5}
    \draw (0.1,\y) -- (-0.1,\y) node[left] {\small $\y$};

  \draw[thick] (-6,0) -- (0,4) -- (0,0) -- (-6,0);

  \draw[->, red, very thick] (-6,0) -- (-3,2) node[right] {$\vec{v}$};

  \fill (-6,0) circle (1.5pt) node[below left]{A};
  \fill (-3,2) circle (1.5pt) node[above right]{B};

\end{tikzpicture}
\caption{Vector $\vec{v}$ from point A to B, shown on a grid with axes for measurement.}
\end{figure}

Now that we have introduced a grid into the diagram, we can use it as a tool for measurement. 
Each square in the grid can be thought of as representing one unit of length.

\medskip

With this in place, we no longer need to guess how long the vector is. 
We can estimate its length by counting how many units it covers.

\medskip

Now we turn to direction. The grid also helps us compare directions in a consistent way. 
The horizontal lines of the grid give us a natural notion of left and right, while the vertical lines give us a notion of up and down.

\medskip

These directions do not come from the vector itself; they are choices that we make in order to describe motion more precisely.

\medskip

Once these reference directions are fixed, we can describe any vector by saying how much it moves in each of these two directions.

\medskip

In other words, instead of thinking of the vector as a single arrow, we can think of it as being built from two simpler motions:
one along the horizontal direction, and one along the vertical direction.

\medskip

Using the grid, we can now read off these motions directly. 
We can see how far the vector moves horizontally, and how far it moves vertically.

\medskip

This allows us to assign numbers to the vector. These numbers tell us how much of each basic direction is needed to reconstruct the original arrow.

\medskip

Thus, while the vector itself is a geometric object, our description of it depends on the choices we make for measuring distance and describing direction.

\medskip

From the grid, we can read off that the vector moves 3 units in the horizontal direction and 2 units in the vertical direction.

\medskip

This raises an important question: why were we able to describe the vector in this way?

\medskip

The reason is that we have chosen two special directions in the plane: one horizontal and one vertical. 
These directions act as our fundamental building blocks. By combining them, we are able to describe any displacement in the plane.

% (rest of document continues exactly in same cleaned style...)

\section{Basis and Vector Spaces}

The two directions we chose, horizontal and vertical, play a special role. 
They allow us to describe every vector in the plane by breaking it into simpler pieces.

We call such a collection of directions a \textbf{basis}.


Very simply, a \textbf{basis} is a set of directions with the property that every vector can be built by combining them in suitable amounts.


\subsection{A Precise Definition of a Basis}

So far, we have been using the horizontal and vertical directions as our basic building blocks. 
We saw that by combining these directions, we were able to construct the vector $\vec{v}$.

We now want to make this idea precise.

To do this, we introduce the notion of a \textbf{vector space}. 
A vector space is a collection of objects, called vectors, together with two operations:

\medskip

Despite the name, a vector space is not necessarily a physical space like the room 
around you or the hill we have been walking on. It is an abstract mathematical structure 
--- a conceptual ``space'' whose elements are called vectors.

\medskip

These vectors need not be arrows or displacements. They could be functions, polynomials, 
matrices, waves, or many other mathematical objects. What matters is not what the objects 
physically are, but that they obey certain rules for combining and scaling them.

\medskip

A vector space therefore consists of a collection of objects, called vectors, together 
with two operations:

\begin{enumerate}
  \item vector addition, which combines two vectors to form another vector, and
  \item scalar multiplication, which scales a vector by a number.
\end{enumerate}

These operations satisfy certain properties (such as associativity, distributivity, and the existence of a zero vector), 
which ensure that vectors behave in a consistent algebraic way.

\medskip

Within a vector space, we are interested in describing vectors in terms of simpler building blocks.

\subsection{Linear Combinations}

Given vectors $\mathbf{v}_1, \mathbf{v}_2, \dots, \mathbf{v}_n$, 
we can form a new vector $\vec{w}$ by taking a linear combination:

\[
\vec{w} = a_1 \mathbf{v}_1 + a_2 \mathbf{v}_2 + \cdots + a_n \mathbf{v}_n,
\]

where $a_1, a_2, \dots, a_n$ are real numbers.

This represents combining the vectors by scaling each one and then adding the results.

\medskip

\subsection{Spanning}

A collection of vectors $\{\mathbf{v}_1, \dots, \mathbf{v}_n\}$ is said to \textbf{span} the vector space 
if every vector in the space can be written as a linear combination of them.

In other words, these vectors are sufficient to build every vector in the space.

\medskip

\subsection{Linear Independence}

However, not all spanning sets are equally useful. 
We would like our building blocks to be as efficient as possible.

A set of vectors $\{\mathbf{v}_1, \dots, \mathbf{v}_n\}$ is said to be \textbf{linearly independent} 
if the only way to write

\[
a_1 \mathbf{v}_1 + a_2 \mathbf{v}_2 + \cdots + a_n \mathbf{v}_n = \mathbf{0}
\]

is by taking

\[
a_1 = a_2 = \cdots = a_n = 0.
\]

This means that none of the vectors can be built from the others. 
Each one contributes something genuinely new.

\medskip

\subsection{Definition of a Basis}

We are now ready to define a basis.

A set of vectors $\{\mathbf{v}_1, \dots, \mathbf{v}_n\}$ is called a \textbf{basis} of a vector space if:

\begin{enumerate}
  \item it spans the space, and
  \item it is linearly independent.
\end{enumerate}

\medskip

This means that a basis is a minimal set of building blocks from which every vector in the space can be constructed.

\medskip

\subsection{Connection to Our Example}

In our diagram, the horizontal direction $\mathbf{e}_1$ and the vertical direction $\mathbf{e}_2$ form a basis of the plane.

Every vector can be written uniquely as a combination of these two directions. 
For example,

\[
\vec{v} = 3\mathbf{e}_1 + 2\mathbf{e}_2.
\]

The numbers $(3,2)$ are called the \textbf{components} of the vector with respect to this basis.
These components depend on the choice of basis, even though the vector itself does not.
\medskip

It is important to note that the choice of basis is not unique. 
We could have chosen a different pair of directions, and we would still be able to describe every vector in the plane. 
However, the numerical components of the vector would then change accordingly.


\subsection{The Hill Example Revisited}

We now return to our hill example and reinterpret it using the language of vector spaces and bases.

\medskip
 
\textbf{What is the vector space here?}

In this example, we are working in the plane of the diagram. 
The vectors we are considering are all possible displacements that can be drawn on this plane.

That is, starting from any point, we can imagine moving in any direction and by any amount. 
All such displacements form our vector space.

This space is two-dimensional, since two independent directions are sufficient to describe any motion within it.

\medskip

\subsection{The Basis in This Example}

When we introduced the grid, we implicitly chose two special directions:

\begin{itemize}
  \item the horizontal direction (along the $x$-axis), and
  \item the vertical direction (along the $y$-axis).
\end{itemize}

We denote these directions by $\mathbf{e}_1$ and $\mathbf{e}_2$, respectively.

These two vectors form a basis of the plane. 
They are linearly independent, and together they span all possible directions in the diagram.

\medskip

\subsection{Writing the Vector in This Basis}

Now consider the vector $\vec{v}$ from point A to point B.

From the grid, we can see that:

\begin{itemize}
  \item the vector moves 3 units in the horizontal direction, and
  \item the vector moves 2 units in the vertical direction.
\end{itemize}

This means that $\vec{v}$ can be written as a linear combination of the basis vectors:

\[
\vec{v} = 3\mathbf{e}_1 + 2\mathbf{e}_2.
\]

This expression tells us exactly how to reconstruct the vector:
we take 3 copies of the horizontal direction and 2 copies of the vertical direction, and add them together.

\medskip

\subsection{A Deeper Observation}

It is important to note that the vector $\vec{v}$ itself is not the pair of numbers $(3,2)$.

Rather, $(3,2)$ is the description of the vector with respect to the chosen basis $\{\mathbf{e}_1, \mathbf{e}_2\}$.

If we had chosen a different pair of directions as our basis, the same vector would be described by a different pair of numbers.

\medskip

Thus, the vector is a geometric object, while its components depend on the basis used to describe it.


\section{Changing the Basis}

So far, we have described vectors using the horizontal and vertical directions as our basis. 
This allowed us to write the vector $\vec{v}$ as

\[
\vec{v} = 3\mathbf{e}_1 + 2\mathbf{e}_2.
\]

However, there is nothing special about these particular directions. 
We are free to choose a different set of directions, as long as they are linearly independent and span the space.

\subsection{A New Choice of Basis}

Suppose instead of choosing purely horizontal and vertical directions, 
we choose directions that are tilted along the hill.

For example, let us define:

\begin{itemize}
  \item $\mathbf{u}_1$ as a direction pointing up along the slope of the hill,
  \item $\mathbf{u}_2$ as another direction that is not parallel to $\mathbf{u}_1$.
\end{itemize}

As long as these two directions are not multiples of each other, they will still form a basis of the plane.

\medskip

\subsection{Describing the Same Vector}

Now we try to describe the same vector $\vec{v}$ using this new basis $\{\mathbf{u}_1, \mathbf{u}_2\}$.

Instead of moving purely horizontally and then vertically, 
we now think of the vector as being built from motions along these new directions.

This means we can write

\[
\vec{v} = a\,\mathbf{u}_1 + b\,\mathbf{u}_2
\]

for some numbers $a$ and $b$.

\medskip

\subsection{What Has Changed?}

The vector $\vec{v}$ has not changed. 
It still represents the same displacement from A to B.

What has changed is only the way we describe it.

Previously, we said the vector was made of 3 units in the horizontal direction and 2 units in the vertical direction.

Now, we describe it in terms of how much it moves along $\mathbf{u}_1$ and $\mathbf{u}_2$.

So the numbers $(3,2)$ will generally be replaced by a different pair of numbers $(a,b)$.

\medskip

\subsection{A Key Insight}

This leads to an important conclusion:

\begin{itemize}
  \item The vector itself is independent of the basis.
  \item The components of the vector depend on the basis.
\end{itemize}

Thus, changing the basis changes the numerical description of the vector, 
but not the geometric object itself.


\subsection{A Concrete Example}

Let us choose a new basis:

$\mathbf{u}_1$ and $\mathbf{u}_2$


We now try to express $\vec{v} = 3\mathbf{e}_1 + 2\mathbf{e}_2$ in terms of $\mathbf{u}_1$ and $\mathbf{u}_2$.


\begin{figure}[H]
 \centering
 \begin{tikzpicture}[scale=0.9]

   % Grid
   \draw[step=1cm, gray!30, thin] (-7,-1) grid (3,5);

   % Axes
   \draw[->, thick] (-7,0) -- (3.5,0) node[right] {$x$};
   \draw[->, thick] (0,-1) -- (0,5.5) node[above] {$y$};

   % Tick marks and labels on x-axis
   \foreach \x in {-6,-5,-4,-3,-2,-1,1,2,3}
     \draw (\x,0.1) -- (\x,-0.1) node[below] {\small $\x$};

   % Tick marks and labels on y-axis
   \foreach \y in {1,2,3,4,5}
     \draw (0.1,\y) -- (-0.1,\y) node[left] {\small $\y$};

   % Hill
   \draw[thick] (-6,0) -- (0,4) -- (0,0) -- (-6,0);

   % Points
   \fill (-6,0) circle (1.5pt) node[below left] {$A$};
   \fill (-3,2) circle (1.5pt) node[above right] {$B$};

   % Vector v
   \draw[->, red, very thick] (-6,0) -- (-3,2) node[above] {$\vec{v}$};

   % Standard basis vectors from A
   \draw[->, violet, thick] (-6,0) -- (-5,0) node[below] {$\mathbf e_1$};
   \draw[->, violet, thick] (-6,0) -- (-6,1) node[left] {$\mathbf e_2$};

   % Rotated basis vectors from A (45 degrees anticlockwise)
   \draw[->, blue, thick] (-6,0) -- (-5.2928932188134,0.7071067811865) node[above right] {$\mathbf u_1$};
   \draw[->, blue!60!black, thick] (-6,0) -- (-6-0.7071,0+0.7071) node[above left] {$\mathbf u_2$};

 \end{tikzpicture}
 \caption{The vector $\vec{v}$ together with the standard basis $\{\mathbf e_1,\mathbf e_2\}$ and the rotated basis $\{\mathbf{u}_1, \mathbf{u}_2\}$ obtained by rotating the standard basis anticlockwise by $45^\circ$.}
 \label{fig:rotated_basis}
\end{figure}

Since the new basis vectors are obtained by rotating the standard basis by $45^\circ$, we have

\[
\mathbf{u}_1 = \frac{1}{\sqrt{2}}(\mathbf{e}_1 + \mathbf{e}_2),
\qquad
\mathbf{u}_2 = \frac{1}{\sqrt{2}}(-\mathbf{e}_1 + \mathbf{e}_2).
\]

We now solve for $\mathbf{e}_1$ and $\mathbf{e}_2$ in terms of $\mathbf{u}_1$ and $\mathbf{u}_2$.

Adding the two equations:

\[
\mathbf{u}_1 + \mathbf{u}_2 = \frac{1}{\sqrt{2}}(2\mathbf{e}_2),
\]

so

\[
\mathbf{e}_2 = \frac{1}{\sqrt{2}}(\mathbf{u}_1 + \mathbf{u}_2).
\]

Subtracting the second equation from the first:

\[
\mathbf{u}_1 - \mathbf{u}_2 = \frac{1}{\sqrt{2}}(2\mathbf{e}_1),
\]

so

\[
\mathbf{e}_1 = \frac{1}{\sqrt{2}}(\mathbf{u}_1 - \mathbf{u}_2).
\]

Now substitute into $\vec{v} = 3\mathbf{e}_1 + 2\mathbf{e}_2$:

\[
\vec{v} = 3\left(\frac{1}{\sqrt{2}}(\mathbf{u}_1 - \mathbf{u}_2)\right)
+ 2\left(\frac{1}{\sqrt{2}}(\mathbf{u}_1 + \mathbf{u}_2)\right).
\]

Factor out $\frac{1}{\sqrt{2}}$:

\[
\vec{v} = \frac{1}{\sqrt{2}} \left[ 3(\mathbf{u}_1 - \mathbf{u}_2) + 2(\mathbf{u}_1 + \mathbf{u}_2) \right].
\]

Expand:

\[
\vec{v} = \frac{1}{\sqrt{2}} \left( 3\mathbf{u}_1 - 3\mathbf{u}_2 + 2\mathbf{u}_1 + 2\mathbf{u}_2 \right).
\]

Simplify:

\[
\vec{v} = \frac{1}{\sqrt{2}} \left( 5\mathbf{u}_1 - \mathbf{u}_2 \right).
\]

Thus, in the rotated basis, the vector is

\[
\vec{v} = \frac{5}{\sqrt{2}}\,\mathbf{u}_1 - \frac{1}{\sqrt{2}}\,\mathbf{u}_2.
\]

This shows that the same vector can have different numerical descriptions depending on the basis used. 
The components are not intrinsic to the vector; they arise from our choice of basis.

    \chapter{Linear Transformations}
    
    \section{What Is \texorpdfstring{$\mathbb{R}^n$}{R to the n}?}
    
    Before defining a linear transformation, we need to clarify the space we are working in.
    
    \medskip
    
    So far, we have been working on the hill, which lives in a two-dimensional world. 
    To describe a point on the hill, we used two numbers: one for horizontal position 
    and one for vertical height. This is what we call $\mathbb{R}^2$.\footnote{The symbol 
$\mathbb{R}$ stands for the set of \textit{real numbers}, which includes positive 
numbers, negative numbers, fractions, irrational numbers such as $\sqrt{2}$ and $\pi$, 
and zero.}.

\medskip

More generally, $\mathbb{R}^n$ is a space where every point is described using $n$ 
numbers. Each number tells you how far you move in one independent direction.

\medskip

\begin{itemize}
  \item In $\mathbb{R}^2$, we need two numbers (like $x$ and $y$) to locate a point on the hill.
  \item In $\mathbb{R}^3$, we need three numbers (for example, left-right, forward-backward, and up-down).
  \item In $\mathbb{R}^n$, we imagine $n$ independent directions, and each vector tells us how far to move along each one.
\end{itemize}

\medskip

Even though we cannot picture higher dimensions, the idea is the same: 
a vector in $\mathbb{R}^n$ is simply a direction and displacement described using $n$ independent measurements.

\medskip

In this chapter, we will continue to use the hill (which lies in $\mathbb{R}^2$) 
for intuition, but everything we say applies equally well in any dimension.

\section{Changing Vectors Themselves}

So far, we have been changing the way we \textit{look} at vectors by choosing different bases. 
The vectors themselves remained the same, but the numbers used to describe them changed.
Before now, we have focused on describing vectors: choosing a basis, assigning coordinates, 
and understanding how those coordinates change when we change our point of view.

\medskip

But in many situations, the vectors themselves are not fixed. The space we are working in 
can change.

\medskip

For example, imagine the hill is made of a flexible material. If the ground stretches 
in the horizontal direction and compresses in the vertical direction, then every point 
on the hill moves. A path that once went from $A$ to $B$ is no longer the same path --- 
it has been deformed along with the hill.

\medskip

In this situation, we are not simply describing the same vector differently. 
The vector itself has changed.

\medskip



We now ask a different question: \textbf{\textit{What if we change the vectors themselves?}}


\medskip

A \textbf{linear transformation} is a function $T: \mathbb{R}^n \to \mathbb{R}^n$ that maps vectors 
to new vectors, while preserving two essential properties:

\begin{itemize}
  \item \textbf{Additivity:} $T(\vec{u} + \vec{v}) = T(\vec{u}) + T(\vec{v})$
  \item \textbf{Scalar homogeneity:} $T(c\vec{v}) = cT(\vec{v})$
\end{itemize}

Geometrically, this means: \textit{lines stay lines, the origin stays fixed, and the grid 
remains evenly spaced} --- it may be stretched, rotated, sheared, or reflected, but it 
never bends or curves.

% ---------------------------------------------------------
\section{Matrices as Representations of Transformations}
% ---------------------------------------------------------

This raises a natural question: how do we describe such a transformation using numbers?

\medskip

We already know that any vector can be written in terms of a basis. For example,

\[
\vec{v} = x\mathbf{e}_1 + y\mathbf{e}_2.
\]

If we apply a linear transformation $T$ to this vector, the two linearity properties give

\[
T(\vec{v}) = T(x\mathbf{e}_1 + y\mathbf{e}_2) = x\,T(\mathbf{e}_1) + y\,T(\mathbf{e}_2).
\]

This is the key idea: \textit{to understand $T$ on every vector, it is enough to know 
what it does to the basis vectors.}

\medskip

Let us write the transformed basis vectors in coordinates:

\[
T(\mathbf{e}_1) = \begin{bmatrix} a \\ c \end{bmatrix}, 
\qquad
T(\mathbf{e}_2) = \begin{bmatrix} b \\ d \end{bmatrix}.
\]

Now consider an arbitrary vector $\vec{v} = \begin{bmatrix} x \\ y \end{bmatrix}$. 
Using the expression above,

\[
T(\vec{v}) = x \begin{bmatrix} a \\ c \end{bmatrix} + y \begin{bmatrix} b \\ d \end{bmatrix}
=
\begin{bmatrix}
a & b \\
c & d
\end{bmatrix}
\begin{bmatrix}
x \\
y
\end{bmatrix}.
\]

The matrix

\[
M =
\begin{bmatrix}
a & b \\
c & d
\end{bmatrix}
\]

is called the \textbf{matrix representation} of the transformation $T$. Each column encodes 
where a basis vector is sent:

In simple terms, this means the following.

\medskip

Take the first basis vector $\mathbf{e}_1 = \begin{bmatrix}1 \\ 0\end{bmatrix}$. 
Apply the transformation $T$ to it. The result is the first column of the matrix:

\[
T(\mathbf{e}_1) = \begin{bmatrix} a \\ c \end{bmatrix}.
\]

\medskip

Similarly, take the second basis vector $\mathbf{e}_2 = \begin{bmatrix}0 \\ 1\end{bmatrix}$. 
Applying $T$ gives the second column:

\[
T(\mathbf{e}_2) = \begin{bmatrix} b \\ d \end{bmatrix}.
\]

\medskip

So each column of the matrix simply tells us what happens to one of the basic directions 
of the space. Once we know how these basic directions move, we can determine how 
\textit{every} vector moves.

\begin{itemize}
  \item The \textbf{first column} is $T(\mathbf{e}_1)$.
  \item The \textbf{second column} is $T(\mathbf{e}_2)$.
\end{itemize}

A matrix is therefore not merely a collection of numbers --- it is a compact encoding of 
how a transformation acts on the entire space.

% ---------------------------------------------------------
\section{A Worked Example}
% ---------------------------------------------------------

Consider a transformation $T$ that stretches space horizontally by a factor of $1.4$ and 
compresses it vertically by a factor of $0.7$. To find its matrix, we apply $T$ to each 
standard basis vector:

\begin{itemize}
  \item $\mathbf{e}_1 = \begin{bmatrix} 1 \\ 0 \end{bmatrix}$ points entirely in the 
  horizontal direction, so it is stretched to 
  $T(\mathbf{e}_1) = \begin{bmatrix} 1.4 \\ 0 \end{bmatrix}$.
  \item $\mathbf{e}_2 = \begin{bmatrix} 0 \\ 1 \end{bmatrix}$ points entirely in the 
  vertical direction, so it is compressed to 
  $T(\mathbf{e}_2) = \begin{bmatrix} 0 \\ 0.7 \end{bmatrix}$.
\end{itemize}

Placing these as columns gives the matrix that completely describes $T$:

\[
M = \begin{bmatrix} 1.4 & 0 \\ 0 & 0.7 \end{bmatrix}.
\]

We can now apply this to any vector. Consider the point $A = (-6, 0)$ and the point 
$B = (-3, 2)$, connected by a vector $\vec{v}$.
Applying $T$ gives:

\[
T(A) = \begin{bmatrix} 1.4 & 0 \\ 0 & 0.7 \end{bmatrix} \begin{bmatrix} -6 \\ 0 \end{bmatrix} 
= \begin{bmatrix} -8.4 \\ 0 \end{bmatrix}, 
\qquad
T(B) = \begin{bmatrix} 1.4 & 0 \\ 0 & 0.7 \end{bmatrix} \begin{bmatrix} -3 \\ 2 \end{bmatrix} 
= \begin{bmatrix} -4.2 \\ 1.4 \end{bmatrix}.
\]

These are the \textbf{absolute coordinates} of the transformed points --- that is, where 
$A$ and $B$ end up when measured against the \textit{original}, undeformed coordinate system.

\medskip

\begin{figure}[H]
\centering
\resizebox{\textwidth}{!}{%
\begin{tikzpicture}[scale=0.8]

% =========================
% LEFT: Original figure
% =========================
\begin{scope}
  \node at (-2.5, 6.2) {\textbf{Original}};
  \draw[step=1cm, gray!30, thin] (-7,-1) grid (2,5);
  \draw[->, thick] (-7.2,0) -- (2.5,0) node[right] {$x$};
  \draw[->, thick] (0,-1.2) -- (0,5.5) node[above] {$y$};
  \foreach \x in {-6,-5,-4,-3,-2,-1,1,2}
    \draw (\x,0.1) -- (\x,-0.1) node[below] {\small $\x$};
  \foreach \y in {1,2,3,4,5}
    \draw (0.1,\y) -- (-0.1,\y) node[left] {\small $\y$};
  \draw[thick] (-6,0) -- (0,4) -- (0,0) -- cycle;
  \draw[->, red, very thick] (-6,0) -- (-3,2) node[above right] {$\vec{v}$};
  \fill (-6,0) circle (1.5pt) node[above left] {$A(-6,0)$};
  \fill (-3,2) circle (1.5pt) node[below right] {$B(-3,2)$};
\end{scope}

% =========================
% Arrow + label between figures
% =========================
\node at (3.2, 2.5) {$T$};
\draw[->, thick] (2.2, 2.2) -- (4.2, 2.2);

% =========================
% RIGHT: Transformed figure
% T = [[1.4, 0], [0, 0.7]]
% (-6,0) -> (-8.4, 0)
% (-3,2) -> (-4.2, 1.4)
% (0,4)  -> (0,  2.8)
% (0,0)  -> (0,  0)
% =========================
\begin{scope}[xshift=14cm]
  \node at (-2.5, 6.2) {\textbf{After $T$: stretch $x$ by 1.4, squeeze $y$ by 0.7}};
  \foreach \x in {-7,-6,-5,-4,-3,-2,-1,0,1,2}
    \draw[gray!30, thin] ({1.4*\x}, {0.7*(-1)}) -- ({1.4*\x}, {0.7*5});
  \foreach \y in {-1,0,1,2,3,4,5}
    \draw[gray!30, thin] ({1.4*(-7)}, {0.7*\y}) -- ({1.4*2}, {0.7*\y});
  \draw[->, thick] ({1.4*(-7)-0.2}, 0) -- ({1.4*2+0.5}, 0) node[right] {$x'$};
  \draw[->, thick] (0, {0.7*(-1)-0.2}) -- (0, {0.7*5+0.5}) node[above] {$y'$};
  \foreach \x in {-7,-6,-5,-4,-3,-2,-1,1,2}
    \draw ({1.4*\x}, 0.1) -- ({1.4*\x}, -0.1) node[below] {\small $\x$};
  \foreach \y in {1,2,3,4,5}
    \draw (0.1, {0.7*\y}) -- (-0.1, {0.7*\y}) node[left] {\small $\y$};
  \draw[thick] (-8.4,0) -- (0,2.8) -- (0,0) -- cycle;
  \draw[->, red, very thick] (-8.4,0) -- (-4.2,1.4) node[above right] {$T(\vec{v})$};
  \fill (-8.4,0) circle (1.5pt) node[above left] {$T(A)(-8.4,\,0)$};
  \fill (-4.2,1.4) circle (1.5pt) node[below right] {$T(B)(-4.2,\,1.4)$};
\end{scope}

\node at (7, -1.8) {$T = \begin{bmatrix} 1.4 & 0 \\ 0 & 0.7 \end{bmatrix}$};

\end{tikzpicture}%
}
\caption{A linear transformation $T$ that stretches space horizontally by 1.4 and compresses 
it vertically by 0.7. The triangle and vector transform consistently --- grid lines, shape, 
and vector all obey the same mapping.}
\label{fig:linear_transform}
\end{figure}

% ---------------------------------------------------------
\section{The Two Ways to Read the Transformed Grid}
% ---------------------------------------------------------

After applying $T$, we are looking at the same geometric picture, but we can read it in two different ways.

\medskip

Now you apply a linear transformation to the grid.

\medskip

\textbf{Reading 1: Using the original grid.}

\medskip

In this interpretation, we are measuring everything using the original grid.

\medskip

The transformation $T$ physically moves every point. In particular,

\[
A = (-6, 0) \quad \longrightarrow \quad T(A) = (-8.4, 0).
\]

\medskip

So relative to the original grid, point $A$ has clearly moved further to the left.

\medskip

\textbf{Reading 2: Using the transformed grid.}

\medskip

In this interpretation, we imagine that the grid itself has been stretched by the transformation. 
We now use this new grid.

\medskip

When we measure $T(A)$ using this stretched grid, we find that it still lines up with the 
same tick mark as before. So we describe it as

\[
(-6, 0).
\]

\medskip

\textbf{Key idea.} The point has moved, but the grid has also changed. Because both changed 
in the same way, the coordinates appear unchanged.

\medskip

\begin{center}
\begin{tabular}{l@{\hspace{2cm}}c}
\toprule
& \textbf{Coordinates} \\
\midrule
$A$ in the original grid        & $(-6,\ 0)$ \\
$T(A)$ in the original grid      & $(-8.4,\ 0)$ \\
$T(A)$ measured using the new grid   & $(-6,\ 0)$ \quad \textit{--- looks the same!} \\
\bottomrule
\end{tabular}
\end{center}

\medskip

\textbf{Why this looks confusing.} In the transformed picture, the spacing of the grid has changed. 
The distance between $0$ and $1$ on the $x$-axis is now larger than before. 

\medskip

So even though $T(A)$ sits on the ``$-6$'' tick mark in the new grid, that tick mark is physically 
farther away than it was in the original grid.

\medskip

\subsection{Analogy: Floor Tiles}

To understand the two readings, imagine the floor of a room covered with square tiles. 
Each tile represents one unit of distance, and we describe positions by counting how many tiles we move.

\medskip

Suppose a point is located 6 tiles to the left of the origin. We describe it as $(-6, 0)$.

\medskip

Now imagine that the floor is stretched in the horizontal direction so that each tile becomes wider. 
The pattern of tiles remains, but the spacing between them increases.

\medskip

There are now two ways to describe the position of the same point:

\begin{itemize}
  \item If we measure using the original floor (before stretching), the point is now farther away. 
  It lies at $(-8.4, 0)$.

  \item If we measure using the stretched tiles, the point still sits on the 6th tile. 
  So we describe it as $(-6, 0)$.
\end{itemize}

\medskip

\textbf{Key idea.} The point has not changed relative to the tile it sits on, but the tiles themselves 
have changed size. Because of this, the same point can have different numerical descriptions.

\medskip

This is exactly what happens in our transformed grid:

\begin{itemize}
  \item The original grid corresponds to the original tiling.
  \item The transformed grid corresponds to the stretched tiling.
\end{itemize}

Both describe the same geometric situation, but they assign different coordinates because 
they use different grids.

\medskip

The diagram shows the \textbf{active transformation} interpretation --- the whole space 
stretches, and $T(A) = (-8.4, 0)$ tells you where $A$ lands in the original, undeformed 
coordinate system. The stretched grid is a visual aid to show \textit{how} space deformed, 
but the labels $T(A)$ and $T(B)$ are always reported back in the original coordinate language. 
This is why linear transformations are powerful: they tell you exactly where every point 
ends up in absolute terms, even as the grid itself warps.

% ---------------------------------------------------------
\section{Linear Transformations vs.\ Change of Basis}
% ---------------------------------------------------------

It is natural to ask: if rotating the grid $90^\circ$ via a linear transform produces the 
same numbers as tilting your head $90^\circ$ (a change of basis), are they not doing the 
same thing?

\medskip

The answer is that the \textit{numerical result} can be identical, but the interpretation 
is fundamentally different. For orthogonal transformations such as rotations, the active 
and passive interpretations are mathematically equivalent --- same matrix, same numbers, 
two different stories. This is in fact the standard way physicists think about it, and they 
have a name for it: the \textbf{active} versus \textbf{passive} interpretation of a 
transformation.

\medskip

The difference becomes visible the moment there is anything \textbf{else} on the table. 
\subsection{Analogy: The Table and Your Head}

To understand the two viewpoints, it helps to think of a simple analogy.

\medskip

Imagine there is a table in a room, and on the table lies a pen pointing toward the door. Let's say the direction the pen is pointing in is your y-direction.
You are standing in the room looking at the table.

\medskip

Now consider two different situations.

\medskip

\textbf{Scenario 1: The table moves.}

\begin{itemize}
  \item The table is physically rotated.
  \item The pen rotates along with the table.
  \item The pen is now not pointing at the door. Instead it is pointing at the window in the room let's say.
\end{itemize}

You will notice that the door is still in your y-direction, but the pen no longer points in that direction and hence no longer toward the door.

This is like a \textbf{linear transformation}: the vector itself is changing.

\medskip

\textbf{Scenario 2: You tilt your head.}

\begin{itemize}
  \item The table does not move at all.
  \item The pen still points toward the door.
  \item But because your viewpoint has changed, the pen appears to point in a different direction relative to your frame of reference.
\end{itemize}

Interestingly, from the perspective of the room, nothing has changed. 
The table is still in the same position, and the pen is still pointing toward the door.

\medskip

However, from your point of view, things appear different. 
Because you have tilted your head, the pen no longer aligns with the directions you are using. 
What once pointed along your y-direction now appears to point in a different direction.

\medskip

Yet, you still agree that the pen is pointing toward the door. 
To make sense of this, you must also recognise that the position of the door in your description has changed. 
In your new way of measuring directions, the door is no longer located at the same coordinates as before.

\medskip

In other words, nothing in the room has changed --- only the way you are describing it has.

This is like a \textbf{change of basis}: the vector stays the same, but the coordinate system changes.

\medskip

\textbf{Key idea.} In one case the pen has physically rotated, and in the other case only your perspective has changed.

\medskip


The real distinguishing question is: \textit{is there an external reference frame that matters?} 
If the table exists in a room with walls and a door, rotating the table changes its 
relationship to the room. A change of basis does not move the table at all --- it only 
changes how you are reading the coordinates. If the table were the entire universe with 
nothing outside it, the two would become indistinguishable.

\medskip

\begin{center}
\begin{tabular}{l@{\hspace{1cm}}c@{\hspace{1cm}}c}
\toprule
& \textbf{Linear Transform} & \textbf{Change of Basis} \\
\midrule
The table   & Physically rotates    & Stays still \\
The vector   & Moves with the table   & Stays still, re-labelled \\
The pen    & Moves with the table   & Stays still, re-labelled \\
Your head   & Stays upright       & Tilts $90^\circ$ \\
The room    & Unchanged --- table moved & Unchanged --- nothing moved \\
\bottomrule
\end{tabular}
\end{center}

\medskip

For a pure rotation in isolation the two are two sides of the same coin. The difference 
only becomes meaningful when there are \textit{multiple objects}, or an \textit{external 
frame of reference}, to relate things to.

% ---------------------------------------------------------
\section{Composing Transformations}
% ---------------------------------------------------------

Suppose we apply one transformation $T_1$, and then apply another transformation $T_2$.

\medskip

Geometrically, this means we first deform the hill in one way, and then deform the result again.

\medskip

The combined transformation is written as

\[
T_2(T_1(\vec{v})).
\]

\medskip

If $T_1$ is represented by the matrix $M_1$, and $T_2$ is represented by the matrix $M_2$, 
then the combined transformation is represented by the product

\[
T_2(T_1(\vec{v})) = M_2 M_1 \vec{v}.
\]

\medskip

\textbf{A concrete example.}

Let

\[
M_1 =
\begin{bmatrix}
1.4 & 0 \\
0 & 0.7
\end{bmatrix}
\]

represent the stretching and squeezing of the hill from earlier.

\medskip

Now let

\[
M_2 =
\begin{bmatrix}
0 & -1 \\
1 & 0
\end{bmatrix}
\]

represent a $90^\circ$ rotation.

\medskip

Applying $M_1$ first stretches the hill. Applying $M_2$ next rotates the already stretched hill.

\medskip

The combined effect is

\[
M_2 M_1 =
\begin{bmatrix}
0 & -1 \\
1 & 0
\end{bmatrix}
\begin{bmatrix}
1.4 & 0 \\
0 & 0.7
\end{bmatrix}
=
\begin{bmatrix}
0 & -0.7 \\
1.4 & 0
\end{bmatrix}.
\]

\medskip

\textbf{Important point.} The order matters. Rotating first and then stretching would produce 
a different result.

\medskip

Thus, matrix multiplication encodes the composition of transformations.

% ---------------------------------------------------------
\section{Changing the Basis of a Transformation}
% ---------------------------------------------------------

We have so far represented transformations using the standard basis 
$\{\mathbf{e}_1, \mathbf{e}_2\}$. But there is nothing special about this choice. 
Any basis will do --- and in fact, for many transformations, a carefully chosen basis 
makes the matrix dramatically simpler to work with. A transformation that looks complicated 
in one basis may become diagonal, or even the identity, in another.

\medskip

\medskip

\textbf{Why is a diagonal matrix special?}

\medskip

A diagonal matrix has the form

\[
\begin{bmatrix}
\lambda_1 & 0 \\
0 & \lambda_2
\end{bmatrix}.
\]

\medskip

This means that each direction is scaled independently:

\begin{itemize}
  \item A vector along $\mathbf{e}_1$ is simply stretched by $\lambda_1$.
  \item A vector along $\mathbf{e}_2$ is simply stretched by $\lambda_2$.
\end{itemize}

\medskip

There is no mixing between directions. Movement in one direction does not produce 
any component in the other.

\medskip

In contrast, a general matrix such as

\[
\begin{bmatrix}
a & b \\
c & d
\end{bmatrix}
\]

mixes directions: moving along $\mathbf{e}_1$ can produce motion in both the 
$\mathbf{e}_1$ and $\mathbf{e}_2$ directions after the transformation.

\medskip

\textbf{Geometric meaning.} A diagonal matrix represents a transformation that 
simply stretches or compresses space along certain preferred directions, without 
rotating or shearing it.

\medskip

These preferred directions are exactly the eigenvectors of the transformation.

\medskip

\textbf{Key idea.} Diagonalising a matrix means choosing a basis aligned with these 
natural directions, so that the transformation acts as pure scaling along each axis.

This raises a central question: if we describe the same transformation using a different 
basis, what does its matrix look like?

\medskip

The answer is the subject of this section. The punchline is:

\[
\boxed{M' = P^{-1} M P}
\]

The transformation itself does not change. Only its matrix representation does, 
depending on which basis you use to describe it.

% ---------------------------------------------------------
\subsection{Deriving \texorpdfstring{$M' = P^{-1}MP$}{M prime equals P inverse M P}}
% ---------------------------------------------------------

Let $\{\mathbf{b}_1, \mathbf{b}_2\}$ be a new basis, and let $P$ be the matrix whose 
columns are these new basis vectors:

\[
P = \begin{bmatrix} \mathbf{b}_1 & \mathbf{b}_2 \end{bmatrix}.
\]
Here, $P$ is not a row matrix. The notation

\[
P = \begin{bmatrix} \mathbf{b}_1 & \mathbf{b}_2 \end{bmatrix}
\]

means that the vectors $\mathbf{b}_1$ and $\mathbf{b}_2$ are placed as \textbf{columns} of the matrix.

\medskip

If
\[
\mathbf{b}_1 = \begin{bmatrix} b_{11} \\ b_{21} \end{bmatrix}, 
\qquad
\mathbf{b}_2 = \begin{bmatrix} b_{12} \\ b_{22} \end{bmatrix},
\]
then the matrix $P$ is

\[
P =
\begin{bmatrix}
b_{11} & b_{12} \\
b_{21} & b_{22}
\end{bmatrix}.
\]

\medskip

So each column of $P$ is one of the basis vectors written in the standard basis.

\medskip

\textbf{Key idea.} The matrix $P$ is formed by stacking the new basis vectors as columns.
$P$ is called the \textbf{change of basis matrix}. Its role is to convert coordinates 
written in the new basis into coordinates in the original standard basis. Concretely, 
if a vector has coordinates $v'$ in the new basis, then its coordinates in the 
standard basis are:

\[
\vec{v} = P\vec{v}'.
\]

Now suppose we want to apply the transformation $M$ to a vector $\vec{v}'$ expressed 
in the new basis, and we want the result expressed in the new basis as well. 
We cannot simply compute $M\vec{v}'$, because $M$ was built to act on standard-basis 
coordinates. We must follow three steps:

\begin{enumerate}
  \item \textbf{Translate into the standard basis.} Convert $\vec{v}'$ into standard 
  coordinates: $\vec{v} = P\vec{v}'$.
  \item \textbf{Apply the transformation.} Apply $M$ in the standard basis: 
  $M\vec{v} = MP\vec{v}'$.
  \item \textbf{Translate back.} Convert the result back into the new basis by 
  applying $P^{-1}$: $P^{-1}MP\vec{v}'$.
\end{enumerate}

\medskip

The composition of these three steps is itself a linear transformation, represented by 
the matrix:

\[
M' = P^{-1}MP.
\]

This is the matrix of the same transformation $M$, but now expressed entirely in the 
new basis. It tells you: given a vector in new-basis coordinates, where does the 
transformation send it, expressed back in new-basis coordinates?

% ---------------------------------------------------------
\subsection{Geometric Interpretation}
% ---------------------------------------------------------

It helps to think of $M' = P^{-1}MP$ as a round trip through three different 
coordinate worlds. The diagram below shows this concretely using the hill.

\medskip
\medskip

\begin{center}
\begin{tikzpicture}[node distance=2.2cm, auto]
  \node (A) at (0,0)  [draw, rounded corners, minimum width=3cm, minimum height=1cm, 
             fill=gray!10] {New basis: $\vec{v}'$};
  \node (B) at (6,0)  [draw, rounded corners, minimum width=3cm, minimum height=1cm, 
             fill=gray!10] {Standard basis: $\vec{v}$};
  \node (C) at (6,-3) [draw, rounded corners, minimum width=3cm, minimum height=1cm, 
             fill=gray!10] {Standard basis: $M\vec{v}$};
  \node (D) at (0,-3) [draw, rounded corners, minimum width=3cm, minimum height=1cm, 
             fill=gray!10] {New basis: $M'\vec{v}'$};
  \draw[->, thick] (A) -- (B) node[midway, above] {$P$};
  \draw[->, thick] (B) -- (C) node[midway, right] {$M$};
  \draw[->, thick] (C) -- (D) node[midway, above] {$P^{-1}$};
  \draw[->, thick, dashed] (A) -- (D) node[midway, left] {$M' = P^{-1}MP$};
\end{tikzpicture}
\end{center}

\medskip
The diagram below shows this concretely using the hill.
\begin{figure}[H]
\centering
\resizebox{\textwidth}{!}{%
\begin{tikzpicture}[scale=0.75]

% ============================================================
% PANEL 1 (top-left): New basis
% ============================================================
\begin{scope}[xshift=0cm, yshift=0cm]

  \node[font=\bfseries] at (-2.5, 7.0) {(1) New basis: $\vec{v}'$};

  \begin{scope}
    \clip (-7.3,-1.3) rectangle (2.3,5.8);
    \foreach \k in {-3,-2,-1,0,1,2,3,4,5} {
      \draw[gray!30, thin]
        ({-8 + \k*(-0.5)}, {-2 + \k*1.0})
       -- ({ 4 + \k*(-0.5)}, { 6 + \k*1.0});
    }
    \foreach \k in {-4,-3,-2,-1,0,1,2,3} {
      \draw[gray!30, thin]
        ({\k*2.0 + (-2)*(-0.5)}, {\k*0.5 + (-2)*1.0})
       -- ({\k*2.0 +  4*(-0.5)}, {\k*0.5 +  4*1.0});
    }
  \end{scope}

  \draw[->, thick] (-7.2,0) -- (2.5,0) node[right] {$x$};
  \draw[->, thick] (0,-1.2) -- (0,5.8) node[above] {$y$};

  \foreach \x in {-6,-5,-4,-3,-2,-1,1,2}
    \draw (\x,0.08) -- (\x,-0.08) node[below] {\small $\x$};
  \foreach \y in {1,2,3,4,5}
    \draw (0.08,\y) -- (-0.08,\y) node[left] {\small $\y$};

  \draw[thick] (-6,0) -- (0,4) -- (0,0) -- cycle;
  \draw[->, red, very thick] (-6,0) -- (-3,2) node[above=3pt] {$\vec{v}'$};

  \fill (-6,0) circle (2pt) node[above left=2pt] {\small $A(-6,0)$};
  \fill (-3,2) circle (2pt) node[right=3pt] {\small $B(-3,2)$};

  \draw[->, blue!80!black, thick] (0,0) -- (1.0,0.9)
    node[right=2pt] {\small $\mathbf{b}_1$};
  \draw[->, blue!80!black, thick] (0,0) -- (-0.5,1.0)
    node[left=2pt] {\small $\mathbf{b}_2$};

  \node[gray!55!black, font=\small\itshape] at (-2.5,-1.1)
    {Tilted grid (new basis)};

\end{scope}

% ============================================================
% Arrow 1->2: P (generous gap between panels)
% ============================================================
\draw[->, very thick] (1.5, 2.8) -- (4.5, 2.8);
\node[font=\Large\bfseries] at (3.0, 3.6) {$P$};

% ============================================================
% PANEL 2 (top-right): Standard basis
% xshift=10cm gives a 10cm gap from origin of panel 1
% Panel 1 runs from x=-7.2 to x=2.5, so rightmost point is at
% absolute x=2.5. Panel 2 starts at absolute x=10, giving
% a clean 7.5cm gap which the arrow fills.
% ============================================================
\begin{scope}[xshift=12cm, yshift=0cm]

  \node[font=\bfseries] at (-2.5, 7.0) {(2) Standard basis: $\vec{v} = P\vec{v}'$};

  \draw[step=1cm, gray!30, thin] (-7.0,-1.0) grid (2.0,5.0);

  \draw[->, thick] (-7.2,0) -- (2.5,0) node[right] {$x$};
  \draw[->, thick] (0,-1.2) -- (0,5.8) node[above] {$y$};

  \foreach \x in {-6,-5,-4,-3,-2,-1,1,2}
    \draw (\x,0.08) -- (\x,-0.08) node[below] {\small $\x$};
  \foreach \y in {1,2,3,4,5}
    \draw (0.08,\y) -- (-0.08,\y) node[left] {\small $\y$};

  \draw[thick] (-6,0) -- (0,4) -- (0,0) -- cycle;
  \draw[->, red, very thick] (-6,0) -- (-3,2) node[above=3pt] {$\vec{v}$};

  \fill (-6,0) circle (2pt) node[above left=2pt] {\small $A(-6,0)$};
  \fill (-3,2) circle (2pt) node[right=3pt] {\small $B(-3,2)$};

  \draw[->, blue!80!black, thick] (0,0) -- (1,0)
    node[below right=2pt] {\small $\mathbf{e}_1$};
  \draw[->, blue!80!black, thick] (0,0) -- (0,1)
    node[above left=2pt] {\small $\mathbf{e}_2$};

  \node[gray!55!black, font=\small\itshape] at (-2.5,-1.1)
    {Standard grid};

\end{scope}

% ============================================================
% Arrow 2->3: M (right side of panel 2, going down)
% Panel 2 is at xshift=10, its right edge is at 10+2.5=12.5
% Arrow sits at x=13.5, safely outside
% ============================================================
\draw[->, very thick] (6.0,-1.5) -- (6.0,-4.0);
\node[font=\Large\bfseries] at (6.4,-2.8) {$M$};

% ============================================================
% PANEL 3 (bottom-right): Stretched standard basis
% Sits at same xshift=10 as panel 2, yshift=-10cm
% ============================================================
\begin{scope}[xshift=13cm, yshift=-10cm]

  \node[font=\bfseries] at (-2.5, 7.0) {(3) Standard basis: $M\vec{v}$};

  \foreach \x in {-7,-6,-5,-4,-3,-2,-1,0,1,2}
    \draw[gray!30, thin] ({1.4*\x},-1.0) -- ({1.4*\x},5.0);
  \foreach \y in {-1,0,1,2,3,4,5}
    \draw[gray!30, thin] ({1.4*(-7)},{0.7*\y}) -- ({1.4*2},{0.7*\y});

  \draw[->, thick] ({1.4*(-7)-0.3},0) -- ({1.4*2+0.8},0) node[right] {$x'$};
  \draw[->, thick] (0,-1.2) -- (0,5.8) node[above] {$y'$};

  \foreach \x in {-6,-5,-4,-3,-2,-1,1,2}
    \draw ({1.4*\x},0.08) -- ({1.4*\x},-0.08) node[below] {\small $\x$};
  \foreach \y in {1,2,3,4}
    \draw (0.08,{0.7*\y}) -- (-0.08,{0.7*\y}) node[left] {\small $\y$};

  \draw[thick] (-8.4,0) -- (0,2.8) -- (0,0) -- cycle;
  \draw[->, red, very thick] (-8.4,0) -- (-4.2,1.4)
    node[above=3pt] {$M\vec{v}$};

  \fill (-8.4,0) circle (2pt) node[above=5pt] {\small $T(A)(-8.4,\ 0)$};
  \fill (-4.2,1.4) circle (2pt) node[below=7pt] {\small $T(B)(-4.2,\ 1.4)$};

  \draw[->, blue!80!black, thick] (0,0) -- (1.4,0)
    node[below right=2pt] {\small $T(\mathbf{e}_1)$};
  \draw[->, blue!80!black, thick] (0,0) -- (0,0.7)
    node[above left=2pt] {\small $T(\mathbf{e}_2)$};

  \node[gray!55!black, font=\small\itshape] at (-2.5,-1.1)
    {Stretched standard grid};

\end{scope}

% ============================================================
% Arrow 3->4: P^{-1} (bottom row, going left)
% Sits at y = -10 + 2.5 = -7.5 (mid-height of bottom panels)
% ============================================================
\draw[->, very thick] (4.5,-7.5) -- (1.5,-7.5);
\node[font=\Large\bfseries] at (3.0,-6.7) {$P^{-1}$};

% ============================================================
% PANEL 4 (bottom-left): Back to new basis
% Same xshift=0 as panel 1, yshift=-10cm
% ============================================================
\begin{scope}[xshift=0cm, yshift=-10cm]

  \node[font=\bfseries] at (-2.5, 7.0)
    {(4) New basis: $M'\vec{v}' = P^{-1}MP\,\vec{v}'$};

  \begin{scope}
    \clip (-7.3,-1.3) rectangle (2.3,5.8);
    \foreach \k in {-3,-2,-1,0,1,2,3,4,5} {
      \draw[gray!30, thin]
        ({-8 + \k*(-0.5)}, {-2 + \k*1.0})
       -- ({ 4 + \k*(-0.5)}, { 6 + \k*1.0});
    }
    \foreach \k in {-4,-3,-2,-1,0,1,2,3} {
      \draw[gray!30, thin]
        ({\k*2.0 + (-2)*(-0.5)}, {\k*0.5 + (-2)*1.0})
       -- ({\k*2.0 +  4*(-0.5)}, {\k*0.5 +  4*1.0});
    }
  \end{scope}

  \draw[->, thick] (-7.2,0) -- (2.5,0) node[right] {$x$};
  \draw[->, thick] (0,-1.2) -- (0,5.8) node[above] {$y$};

  \foreach \x in {-6,-5,-4,-3,-2,-1,1,2}
    \draw (\x,0.08) -- (\x,-0.08) node[below] {\small $\x$};
  \foreach \y in {1,2,3,4,5}
    \draw (0.08,\y) -- (-0.08,\y) node[left] {\small $\y$};

  \draw[thick] (-8.4,0) -- (0,2.8) -- (0,0) -- cycle;
  \draw[->, red, very thick] (-8.4,0) -- (-4.2,1.4)
    node[above=3pt] {$M'\vec{v}'$};

  \fill (-8.4,0) circle (2pt) node[above left=2pt] {\small $T(A)$};
  \fill (-4.2,1.4) circle (2pt) node[right=3pt] {\small $T(B)$};

  \draw[->, blue!80!black, thick] (0,0) -- (2.0,0.5)
    node[right=2pt] {\small $\mathbf{b}_1$};
  \draw[->, blue!80!black, thick] (0,0) -- (-0.5,1.0)
    node[left=2pt] {\small $\mathbf{b}_2$};

  \node[gray!55!black, font=\small\itshape] at (-2.5,-1.1)
    {Tilted grid (new basis)};

\end{scope}

% ============================================================
% Dashed shortcut arrow: runs down the left spine
% Starts above panel 1, ends below panel 4
% At x=-8.5, well outside the clip region of panels 1 and 4
% ============================================================
\draw[->, dashed, very thick, red!70!black]
  (-8.5, 1.5) -- (-8.5, -6.5);
\node[red!70!black, font=\bfseries, rotate=90] at (-9.5, -1.5)
  {$M' = P^{-1}MP$};

\end{tikzpicture}%
}
\caption{The round trip of $M' = P^{-1}MP$. Starting from $\vec{v}'$ in the new basis
\textbf{(1)}, we convert to standard coordinates via $P$ \textbf{(2)}, apply $M$
\textbf{(3)}, then convert back via $P^{-1}$ \textbf{(4)}. The dashed arrow is the
direct path: $M'$ performs all three steps in a single matrix.}
\label{fig:change_of_basis_roundtrip}
\end{figure}

\medskip

The dashed arrow is what $M'$ does directly. The three solid arrows show the long way 
around: translate into the standard basis, apply the transformation, translate back. 
Both paths start and end in the new basis, and both give exactly the same result.


\medskip

Geometrically, think of it this way. Suppose your new basis vectors $\mathbf{b}_1$ 
and $\mathbf{b}_2$ define a tilted, rescaled grid --- like the stretched grid from 
our earlier example. A vector described in this grid has coordinates that mean 
something different from the same numbers in the standard grid.

\medskip

$P$ is the act of \textit{stepping out} of the tilted grid and into the standard room, 
so that $M$ can act on it in the language it understands. $P^{-1}$ is the act of 
\textit{stepping back} into the tilted grid to read off the result. The transformation 
$M'$ is what an observer who lives entirely inside the tilted grid would write down 
to describe exactly the same physical transformation.

% ---------------------------------------------------------

\subsection{Connection to the Hill}
% ---------------------------------------------------------

Return to the hill example. The standard basis $\{\mathbf{e}_1, \mathbf{e}_2\}$ 
describes vectors using the original upright grid. Now suppose we introduce a new 
basis aligned with the hill itself --- for example, one basis vector running along 
the base of the hill and another pointing up the slope.

\medskip

To see why this might be useful, imagine a steady wind blowing across the hill, 
flowing exactly along the slope.

\medskip

If we describe this wind using the standard horizontal and vertical directions, 
its motion appears split between the two --- part of it looks horizontal, and part 
of it looks vertical. The description feels unnecessarily complicated.

\medskip

However, if we choose a basis aligned with the hill --- one direction along the slope 
and one perpendicular to it --- the situation becomes much simpler. The wind now points 
entirely along a single direction in this new basis, with no mixing between components.

\medskip

This is the advantage of choosing a good basis: it aligns our description with the 
natural directions of the problem.

\medskip

In this new basis, $P$ converts hill-aligned coordinates into standard coordinates. 
$M$ applies the stretch-and-squeeze transformation in the standard grid. $P^{-1}$ 
reads the result back in hill-aligned coordinates.

\medskip

The hill has not moved. The transformation has not changed. But the matrix $M'$ 
describes how the transformation looks to someone who measures everything along the 
slope of the hill rather than along the horizontal and vertical axes.

\medskip

\textbf{Key idea.} A good choice of basis reveals the underlying simplicity of a 
transformation by aligning with its natural directions.

\medskip

This is the deeper meaning of $M' = P^{-1}MP$: \textit{same transformation, different 
language}. The physics is invariant; only the coordinate description changes.

% ---------------------------------------------------------
\subsection{Why This Matters}
% ---------------------------------------------------------

The formula $M' = P^{-1}MP$ means we are free to 
choose whatever basis makes a problem easiest.

\medskip


For many transformations, there exists a special basis --- the basis of 
\textbf{eigenvectors} --- in which $M'$ becomes a diagonal matrix:

\[
M' = \begin{bmatrix} \lambda_1 & 0 \\ 0 & \lambda_2 \end{bmatrix}.
\]

In this basis, the transformation simply scales each basis vector independently by 
$\lambda_1$ and $\lambda_2$, with no mixing between directions. Computations that 
would require repeated matrix multiplications in the standard basis reduce to 
trivial scalar arithmetic in the eigenbasis. This is the idea behind 
\textbf{diagonalisation}, which we develop now.

% ---------------------------------------------------------
\chapter{Eigenvectors and Natural Directions}

\section{Eigenvectors and Eigenvalues}
% ---------------------------------------------------------

We have seen that a linear transformation $T$ moves vectors around in space. Most 
vectors get rotated, stretched, and mixed together --- their direction changes after 
the transformation. But some special vectors are different.

\medskip

A non-zero vector $\vec{v}$ is called an \textbf{eigenvector} of a transformation 
$T$ if applying $T$ does not change its direction --- it only scales it:

\[
T(\vec{v}) = \lambda \vec{v}.
\]

The scalar $\lambda$ is called the \textbf{eigenvalue} associated with $\vec{v}$. 
It tells you by how much the vector is stretched ($\lambda > 1$), compressed 
($0 < \lambda < 1$), reversed ($\lambda < 0$), or left unchanged ($\lambda = 1$).

\medskip

\textbf{Geometrically.} An eigenvector is a non-zero vector that lies along a line through the origin 
which the transformation maps back onto the same line. The transformation may stretch 
or flip the vector along that line, but it never rotates it off it.

\medskip

In terms of the matrix $M$ representing $T$, the eigenvector condition becomes:

\[
M\vec{v} = \lambda\vec{v}.
\]

Rearranging:

\[
M\vec{v} - \lambda\vec{v} = \vec{0}
\implies
(M - \lambda I)\vec{v} = \vec{0}.
\]
We introduce $I$ so that $\lambda$ can act on the entire vector $\vec{v}$ as a matrix, allowing us to combine both terms into a single matrix expression.

\medskip

For a non-zero solution $\vec{v}$ to exist, the matrix $(M - \lambda I)$ must be 
singular --- that is, it must have zero determinant:

\[
\det(M - \lambda I) = 0.
\]

% ---------------------------------------------------------
\section{The Determinant}
% ---------------------------------------------------------

We have arrived at the condition

\[
\det(M - \lambda I) = 0.
\]

\subsection{An Intuitive Picture: Unit Cells and Scaling}

The determinant is one of those rare mathematical concepts where the geometric intuition is often more illuminating than the algebraic formula.

\medskip

At its core, the determinant tells you the \textbf{volume scaling factor} of a linear transformation. 
If we take a region of space with volume $1$ and apply a transformation with matrix $A$, 
the volume of the resulting region is exactly

\[
\text{Volume} = |\det(A)|.
\]

\medskip

\subsection{The Change in ``Unit Cells''}

One way to visualize this is to think of space as being built from tiny unit blocks 
(a square in 2D, a cube in 3D).

\medskip

When we apply a transformation:

\begin{itemize}
  \item The matrix moves the edges of this unit block.
  \item The block becomes a parallelogram (in 2D) or a parallelepiped (in 3D).
  \item The determinant measures the area or volume of this new shape.
\end{itemize}

\medskip

For example:

\begin{itemize}
  \item If $\det(A) = 2$, the transformation doubles all volumes.
  \item If $\det(A) = 0$, the transformation squashes space into a lower dimension, 
  collapsing volume to zero.
\end{itemize}

\begin{figure}[htbp]
\centering

% ---------------------------------------------------------
% 2D Diagram
% ---------------------------------------------------------
\begin{minipage}{0.45\textwidth}
\centering
\begin{tikzpicture}[scale=1]

% Original square
\draw[thick, blue] (0,0) rectangle (2,2);

% Transformed parallelogram
\draw[thick, red]
(4,0) -- (6,0.5) -- (7,2.5) -- (5,2) -- cycle;

% Labels
\node[blue] at (1,-0.5) {Unit square};
\node[red] at (5.5,-0.5) {Transformed parallelogram};

% Arrows
\draw[->, thick] (2.5,1) -- (3.5,1)
node[midway, above] {$A$};

\end{tikzpicture}

\caption*{In $\mathbb{R}^2$, the determinant measures how the area of the unit square changes under the transformation.}
\end{minipage}
\hfill
% ---------------------------------------------------------
% 3D Diagram
% ---------------------------------------------------------
\begin{minipage}{0.45\textwidth}
\centering
\begin{tikzpicture}[scale=0.9]

% Original cube
\draw[thick, blue]
(0,0) -- (2,0) -- (2,2) -- (0,2) -- cycle;

\draw[thick, blue]
(0.7,0.7) -- (2.7,0.7) -- (2.7,2.7) -- (0.7,2.7) -- cycle;

\draw[thick, blue]
(0,0) -- (0.7,0.7);

\draw[thick, blue]
(2,0) -- (2.7,0.7);

\draw[thick, blue]
(2,2) -- (2.7,2.7);

\draw[thick, blue]
(0,2) -- (0.7,2.7);

% Transformed parallelepiped
\begin{scope}[xshift=5cm]

\draw[thick, red]
(0,0) -- (2.2,0.5) -- (2.7,2.2) -- (0.5,1.7) -- cycle;

\draw[thick, red]
(0.8,0.8) -- (3,1.3) -- (3.5,3) -- (1.3,2.5) -- cycle;

\draw[thick, red]
(0,0) -- (0.8,0.8);

\draw[thick, red]
(2.2,0.5) -- (3,1.3);

\draw[thick, red]
(2.7,2.2) -- (3.5,3);

\draw[thick, red]
(0.5,1.7) -- (1.3,2.5);

\end{scope}

% Labels
\node[blue] at (1,-0.7) {Unit cube};
\node[red] at (7,-0.7) {Transformed parallelepiped};

% Arrow
\draw[->, thick] (3.2,1.2) -- (4.2,1.2)
node[midway, above] {$A$};

\end{tikzpicture}

\caption*{In $\mathbb{R}^3$, the determinant measures how the volume of the unit cube changes under the transformation.}
\end{minipage}

\caption{A linear transformation reshapes unit cells of space. The determinant measures the scaling of area and volume.}
\label{fig:determinant_geometry}

\end{figure}

\medskip

\subsection{Connection to Eigenvalues}

This picture connects directly to eigenvalues.

\medskip

For many matrices, the determinant can be written as the product of the eigenvalues:

\[
\det(A) = \lambda_1 \lambda_2 \cdots \lambda_n.
\]

\medskip

Each eigenvalue $\lambda_i$ represents the scaling factor along a particular direction 
(the corresponding eigenvector).

\medskip

For example, if one direction is stretched by a factor of $3$ and another by a factor of $2$, 
then the total area scales by

\[
3 \times 2 = 6.
\]

\medskip

The determinant combines these directional scalings into a single number that describes 
the overall change in volume.

\medskip

\subsection{Stretching, Compressing, and Flipping}

Because the determinant is a product of these directional scaling factors:

\begin{itemize}
  \item If $|\det(A)| > 1$, the transformation expands space.
  \item If $0 < |\det(A)| < 1$, it compresses space.
  \item If $\det(A) < 0$, the transformation also reverses orientation (it includes a reflection).
\end{itemize}

\medskip

The negative sign does \textit{not} mean that area or volume itself becomes physically 
negative. Area and volume are still positive quantities. Instead, the sign keeps track 
of the \textbf{orientation} of the transformed space.

\medskip

To understand this, imagine placing two basis vectors $\mathbf{e}_1$ and $\mathbf{e}_2$ 
tail-to-tail. Moving from $\mathbf{e}_1$ to $\mathbf{e}_2$ in the counterclockwise direction 
defines one orientation of the plane. A transformation with positive determinant preserves 
this orientation. A transformation with negative determinant flips it, changing a 
counterclockwise orientation into a clockwise one.

\medskip

In higher dimensions, this idea generalises to the ``handedness'' of space. For example, 
a right-handed coordinate system may become left-handed after a transformation with 
negative determinant.

\medskip

So the determinant contains two pieces of information simultaneously:

\begin{itemize}
  \item its magnitude tells us how much area or volume scales,
  \item its sign tells us whether orientation is preserved or reversed.
\end{itemize}

\begin{figure}[htbp]
\centering

% ---------------------------------------------------------
% Positive determinant
% ---------------------------------------------------------
\begin{minipage}{0.45\textwidth}
\centering
\begin{tikzpicture}[scale=1.1]

% Axes
\draw[->] (-0.5,0) -- (3,0) node[right] {$x$};
\draw[->] (0,-0.5) -- (0,3) node[above] {$y$};

% Basis vectors
\draw[->, thick, blue] (0,0) -- (2,0) node[below] {$\mathbf{e}_1$};
\draw[->, thick, red] (0,0) -- (0,2) node[left] {$\mathbf{e}_2$};

% Orientation arrow
\draw[->, thick] (1,0.3) arc[start angle=0,end angle=90,radius=0.8];

\node at (1,-1) {$\det(A) > 0$};
\node at (1,-1.5) {Orientation preserved};

\end{tikzpicture}
\end{minipage}
\hfill
% ---------------------------------------------------------
% Negative determinant
% ---------------------------------------------------------
\begin{minipage}{0.45\textwidth}
\centering
\begin{tikzpicture}[scale=1.1]

% Axes
\draw[->] (-0.5,0) -- (3,0) node[right] {$x$};
\draw[->] (0,-2.5) -- (0,1) node[above] {$y$};

% Basis vectors
\draw[->, thick, blue] (0,0) -- (2,0) node[below] {$A\mathbf{e}_1$};
\draw[->, thick, red] (0,0) -- (0,-2) node[left] {$A\mathbf{e}_2$};

% Orientation arrow
\draw[->, thick] (0.3,-1) arc[start angle=90,end angle=0,radius=0.8];

\node at (1,-3) {$\det(A) < 0$};
\node at (1,-3.5) {Orientation reversed};

\end{tikzpicture}
\end{minipage}

\caption{
A positive determinant preserves orientation, while a negative determinant reverses it. 
In the second figure, the transformation flips the plane, changing the handedness of the basis.
}
\label{fig:orientation_flip}

\end{figure}

\medskip

\subsection{Why the Determinant Determines Invertibility}

This also explains why the determinant tells us whether a transformation is invertible.

\medskip

If $\det(A) = 0$, then at least one eigenvalue is zero. This means that some direction 
has been completely collapsed.

\medskip

Geometrically, space has been flattened — for example, a three-dimensional object may be 
reduced to a plane.

\medskip

Once this information is lost, it cannot be recovered. This is why such a transformation 
has no inverse.

\subsection{Why Does the Determinant Have to Be Zero?}

We arrived at the equation

\[
(M - \lambda I)\vec{v} = \vec{0}.
\]

We are looking for a \textbf{non-zero} vector $\vec{v}$ that satisfies this equation.

\medskip

Now, think about what this equation means. It says that the matrix $(M - \lambda I)$ 
takes the vector $\vec{v}$ and sends it to the zero vector.

\medskip

There are two possibilities:

\begin{itemize}
  \item The only solution is $\vec{v} = \vec{0}$.
  \item There exists a non-zero vector $\vec{v} \neq \vec{0}$ such that $(M - \lambda I)\vec{v} = \vec{0}$.
\end{itemize}

\medskip

We are interested in the second case, since eigenvectors must be non-zero.

\medskip

\subsection{Invertibility and Solutions}

If the matrix $(M - \lambda I)$ is \textbf{invertible}, then we can multiply both sides of the equation by its inverse:

\[
(M - \lambda I)^{-1}(M - \lambda I)\vec{v} = (M - \lambda I)^{-1}\vec{0}.
\]

This simplifies to

\[
\vec{v} = \vec{0}.
\]

\medskip

This shows that if $(M - \lambda I)$ is invertible, the only solution is the trivial one. 
So no non-zero eigenvector can exist in this case.

\medskip

\subsection{The Key Conclusion}

For a non-zero solution $\vec{v}$ to exist, the matrix $(M - \lambda I)$ must \textbf{not} be invertible.

\medskip

From our earlier discussion, a matrix is not invertible exactly when its determinant is zero. Therefore,

\[
\det(M - \lambda I) = 0.
\]

\medskip

\subsection{Geometric Interpretation}

Geometrically, this means that the transformation $(M - \lambda I)$ collapses space in at least one direction.

\medskip

In other words, there exists a direction $\vec{v}$ that gets mapped to zero. 
This direction is precisely the eigenvector corresponding to the eigenvalue $\lambda$.

% ---------------------------------------------------------
\section{A Worked Example}
% ---------------------------------------------------------

Return to our transformation matrix:

\[
M = \begin{bmatrix} 1.4 & 0 \\ 0 & 0.7 \end{bmatrix}.
\]

The characteristic equation is:

\[
\det(M - \lambda I) 
= \det\begin{bmatrix} 1.4 - \lambda & 0 \\ 0 & 0.7 - \lambda \end{bmatrix}
= (1.4 - \lambda)(0.7 - \lambda) = 0.
\]

The eigenvalues are therefore $\lambda_1 = 1.4$ and $\lambda_2 = 0.7$.

\medskip

For $\lambda_1 = 1.4$:

\[
(M - 1.4I)\vec{v} = 
\begin{bmatrix} 0 & 0 \\ 0 & -0.7 \end{bmatrix}
\begin{bmatrix} v_1 \\ v_2 \end{bmatrix} = \vec{0}
\implies v_2 = 0.
\]

The eigenvector is any scalar multiple of $\begin{bmatrix} 1 \\ 0 \end{bmatrix}$ 
--- the horizontal axis. This makes geometric sense: the transformation stretches 
horizontally, so any purely horizontal vector stays horizontal.

\medskip

For $\lambda_2 = 0.7$:

\[
(M - 0.7I)\vec{v} = 
\begin{bmatrix} 0.7 & 0 \\ 0 & 0 \end{bmatrix}
\begin{bmatrix} v_1 \\ v_2 \end{bmatrix} = \vec{0}
\implies v_1 = 0.
\]

The eigenvector is any scalar multiple of $\begin{bmatrix} 0 \\ 1 \end{bmatrix}$ 
--- the vertical axis. Again this is geometrically clear: the transformation 
compresses vertically, so any purely vertical vector stays vertical.

\medskip

For a diagonal matrix like this one, the eigenvectors are always the standard basis 
vectors and the eigenvalues are always the diagonal entries. The interesting case 
arises when $M$ is not diagonal --- we must solve the characteristic equation to 
find directions that the transformation preserves.

% ---------------------------------------------------------
\section{The Eigenbasis}
% ---------------------------------------------------------

If a matrix $M$ has $n$ linearly independent eigenvectors, we can collect them into 
a basis for the whole space. This is called the \textbf{eigenbasis}.

\medskip

Let the eigenvectors be $\vec{v}_1, \vec{v}_2, \ldots, \vec{v}_n$ with corresponding 
eigenvalues $\lambda_1, \lambda_2, \ldots, \lambda_n$. Form the matrix $P$ whose 
columns are these eigenvectors:

\[
P = \begin{bmatrix} \vec{v}_1 & \vec{v}_2 & \cdots & \vec{v}_n \end{bmatrix}.
\]

In this basis, the transformation has a special form. Since $M\vec{v}_i = \lambda_i 
\vec{v}_i$ for each $i$, we have:

\[
MP = \begin{bmatrix} M\vec{v}_1 & M\vec{v}_2 & \cdots & M\vec{v}_n \end{bmatrix}
  = \begin{bmatrix} \lambda_1\vec{v}_1 & \lambda_2\vec{v}_2 & \cdots & \lambda_n\vec{v}_n \end{bmatrix}
  = P \begin{bmatrix} \lambda_1 & & \\ & \ddots & \\ & & \lambda_n \end{bmatrix}.
\]

\medskip

This is the key algebraic fact: in the eigenbasis, there is no mixing between 
directions. Each basis vector is simply scaled by its own eigenvalue, independently 
of all the others. The transformation has been completely disentangled.

\medskip

\textbf{Geometrically.} Think back to the hill. In the standard basis, our 
stretch-and-squeeze transformation acts on $x$ and $y$ coordinates simultaneously. 
If we worked with a rotated or tilted basis whose vectors were not aligned with the 
transformation's natural axes, the matrix would have off-diagonal entries --- 
the $x$ and $y$ components would mix. Moving to the eigenbasis is like rotating 
your coordinate system to align with those natural axes, so the mixing disappears 
and each direction behaves independently.

% ---------------------------------------------------------
\section{Diagonalisation}
% ---------------------------------------------------------

We now have all the ingredients. From the section on change of basis, we know that 
if $P$ converts from a new basis to the standard basis, then the matrix of the same 
transformation in the new basis is:

\[
M' = P^{-1}MP.
\]

From the eigenbasis discussion, we know that if $P$ is the matrix of eigenvectors, 
then $MP = P\Lambda$ where $\Lambda$ is the diagonal matrix of eigenvalues. 
Multiplying both sides on the left by $P^{-1}$:

\[
P^{-1}MP = P^{-1}P\Lambda = \Lambda.
\]

This gives us the \textbf{diagonalisation theorem}: if $M$ has $n$ linearly 
independent eigenvectors, then:

\[
\boxed{P^{-1}MP = \Lambda = 
\begin{bmatrix} \lambda_1 & & \\ & \ddots & \\ & & \lambda_n \end{bmatrix}}
\]

where $P$ is the matrix of eigenvectors placed as columns, and $\Lambda$ is the 
diagonal matrix of the corresponding eigenvalues.

\medskip

Equivalently, we can write this as:

\[
M = P \Lambda P^{-1}.
\]

This is the \textbf{eigendecomposition} of $M$. It says that any diagonalisable 
matrix can be factored into three interpretable pieces:

\begin{itemize}
  \item $P^{-1}$ --- change coordinates into the eigenbasis.
  \item $\Lambda$ --- apply the transformation (pure scaling, no mixing).
  \item $P$ --- change coordinates back to the standard basis.
\end{itemize}

This is precisely the round trip $M' = P^{-1}MP$ from before, except now we have 
chosen $P$ optimally --- so that the middle step is as simple as possible.

% ---------------------------------------------------------
\section{Why Diagonalisation Matters}
% ---------------------------------------------------------

\textbf{Matrix powers:} Suppose we want to compute $M^k$ --- applying the 
transformation $k$ times. In the standard basis this requires $k$ matrix 
multiplications. Using the eigendecomposition:

\[
M^k = (P\Lambda P^{-1})^k = P \Lambda^k P^{-1},
\]

because the intermediate $P^{-1}P = I$ terms cancel. And $\Lambda^k$ is trivial 
--- you simply raise each diagonal entry to the power $k$:

\[
\Lambda^k = \begin{bmatrix} \lambda_1^k & & \\ & \ddots & \\ & & \lambda_n^k \end{bmatrix}.
\]

This reduces an exponentially expensive computation to a single set of scalar 
exponentiations.

\medskip

\textbf{Geometric interpretation.} In the eigenbasis, repeated application of $M$ 
simply scales each eigendirection independently. Directions with $|\lambda| > 1$ 
expand, directions with $|\lambda| < 1$ contract, and directions with $\lambda = 1$ 
are fixed. The long-run behaviour of the transformation is completely determined by 
which eigenvalue is largest in magnitude --- this is the \textbf{dominant eigenvalue}.

\medskip

\textbf{Connection to the hill.}

\medskip

Think of the transformation as repeatedly reshaping the hill. Each time we apply $M$, 
the hill is stretched in some directions and compressed in others.

\medskip

If we describe this process using the standard horizontal and vertical directions, 
each step mixes the directions. After several applications, it becomes difficult to 
predict how the hill will evolve.

\medskip

However, in the eigenbasis, each direction behaves independently. One direction 
is scaled by a factor $\lambda_1$, and another by $\lambda_2$.

\medskip

Applying the transformation $k$ times simply multiplies these effects:

\begin{itemize}
  \item the first direction is scaled by $\lambda_1^k$,
  \item the second direction is scaled by $\lambda_2^k$.
\end{itemize}

\medskip

So instead of tracking a complicated sequence of deformations, we can understand 
the long-term behaviour of the hill by following how it scales along these special directions.

\medskip

\textbf{Key idea.} In the right basis, repeated transformations reduce to independent scaling along each direction.

\medskip

In our example,
\[
M = \begin{bmatrix} 1.4 & 0 \\ 0 & 0.7 \end{bmatrix}
\]
is already diagonal, so it is already expressed in its eigenbasis. The eigenvectors 
are $\mathbf{e}_1$ and $\mathbf{e}_2$, and the eigenvalues $1.4$ and $0.7$ tell us 
exactly how the hill changes with each application: it stretches horizontally and 
compresses vertically.

\medskip

After $k$ applications,
\[
M^k = \begin{bmatrix} 1.4^k & 0 \\ 0 & 0.7^k \end{bmatrix}.
\]

\medskip

As $k \to \infty$, the horizontal direction grows without bound while the vertical 
direction shrinks to zero. Geometrically, the hill becomes increasingly flattened, 
approaching a horizontal line.

\medskip

This long-term behaviour is immediately visible from the eigenvalues, and would be 
far harder to extract from repeated matrix multiplication in a non-diagonal form.

% ---------------------------------------------------------
\section{When Diagonalisation Fails}
% ---------------------------------------------------------

Not every matrix is diagonalisable. Diagonalisation requires $n$ linearly 
independent eigenvectors. This can fail in two ways:

\begin{itemize}
  \item \textbf{Repeated eigenvalues with too few eigenvectors.} A matrix may 
have a repeated eigenvalue but not enough independent eigenvectors to span the 
space. Such matrices are called \textbf{defective}.

\medskip

For example, consider the matrix

\[
A =
\begin{bmatrix}
1 & 1 \\
0 & 1
\end{bmatrix}.
\]

Its characteristic equation is

\[
\det(A - \lambda I)
=
\begin{vmatrix}
1-\lambda & 1 \\
0 & 1-\lambda
\end{vmatrix}
=
(1-\lambda)^2.
\]

So the matrix has a repeated eigenvalue:
\[
\lambda = 1.
\]

\medskip

To find the eigenvectors, we solve

\[
(A-I)\vec{v} = \vec{0}.
\]

But

\[
A-I =
\begin{bmatrix}
0 & 1 \\
0 & 0
\end{bmatrix},
\]

which gives

\[
\begin{bmatrix}
0 & 1 \\
0 & 0
\end{bmatrix}
\begin{bmatrix}
x \\
y
\end{bmatrix}
=
\begin{bmatrix}
0 \\
0
\end{bmatrix}.
\]

This implies
\[
y = 0,
\]
while $x$ is free.

\medskip

So every eigenvector has the form

\[
\vec{v} =
\begin{bmatrix}
x \\
0
\end{bmatrix},
\]

meaning there is only \textbf{one independent eigendirection}, even though the matrix acts on a two-dimensional space.

\medskip

Because we cannot find two linearly independent eigenvectors, the matrix cannot be diagonalised.
  \item \textbf{Complex eigenvalues.} A real matrix may have complex eigenvalues 
  (occurring in conjugate pairs). In this case the matrix is not diagonalisable 
  over $\mathbb{R}$, but it is diagonalisable over $\mathbb{C}$, and can be 
  reduced to a real \textbf{block diagonal} form involving $2\times 2$ rotation-scaling 
  blocks.
\end{itemize}

\medskip

% ---------------------------------------------------------
% ---------------------------------------------------------
\part{Measurement and Geometry}
\vspace*{0.35\textheight}

\begin{quote}
\centering
``I must be growing small again.''

\medskip

--- Lewis Carroll, \textit{Alice's Adventures in Wonderland}
\end{quote}



Alice's journey is full of changes in size. She grows, shrinks, reaches for doors, 
measures herself against tables, and slowly discovers that size only makes sense when 
there is something to compare it with.

\medskip

This part of the book begins with the same idea. After studying motion and coordinates, 
we now ask how geometry is measured. What does it mean for a vector to have length? 
What does it mean for two directions to be perpendicular? How can one mathematical object 
measure another?

\medskip

To answer these questions, we introduce covectors, dual spaces, and inner products. These 
ideas allow us to turn motion into numbers, compare directions, measure lengths and 
angles, and understand how geometry becomes something more than a picture.

\medskip

This is the next step down the rabbit hole: from describing motion to measuring the 
geometry beneath it.

\chapter{Covectors and Dual Spaces}

\section{What Is a Covector?}
% ---------------------------------------------------------

Throughout this chapter, we have worked with vectors --- objects that represent 
directions and displacements in space. We now introduce a complementary object.

\medskip

Before we do that, we should clarify something important. So far, we have been thinking 
of a vector as an arrow with a direction and a magnitude. This is a useful geometric picture, 
but it is not the most general definition.

\medskip

In mathematics, a vector is simply an element of a \textit{vector space}. It does not have 
to be an arrow. For example, vectors can also be:
\begin{itemize}
  \item polynomials,
  \item functions,
  \item sequences of numbers.
\end{itemize}

\medskip

What makes these objects ``vectors'' is not their shape, but the fact that they can be 
added together and scaled in a consistent way.

\medskip

The arrow picture is just one way of visualising vectors --- one that is especially useful 
for building geometric intuition, as we have done with the hill.

\medskip

\medskip

It is important to note that a covector is itself a kind of vector --- just not in the 
same space as the vectors we have been working with.

\medskip

While vectors represent directions and displacements, covectors belong to a different 
vector space: the space of all linear functions that act on vectors. This space is 
called the \textbf{dual space}.

\medskip

So even though a covector is not an arrow and does not represent a direction, it still 
behaves like a vector in the sense that we can add covectors together and scale them 
by numbers.

\medskip

\textbf{Key idea.} If a vector tells you how to move, a covector tells you something about that movement in the form of a number.

\medskip

To make this precise, we represent a covector as a function that takes a vector as input 
and returns a number:

\[
f: \mathbb{R}^n \to \mathbb{R}, \qquad \vec{v} \mapsto f(\vec{v}) \in \mathbb{R}.
\]

\medskip

Since a covector measures vectors in a consistent way, this function must be linear. 
That is, for any vectors $\vec{u}, \vec{v}$ and scalar $c$:

\[
f(\vec{u} + \vec{v}) = f(\vec{u}) + f(\vec{v}), \qquad 
f(c\vec{v}) = c\,f(\vec{v}).
\]

A covector does not point anywhere. It does not have a direction in the way a vector 
does. Instead, it \textit{measures} vectors --- it takes a vector and extracts a 
number from it.

% ---------------------------------------------------------
\section{A Concrete Example}
% ---------------------------------------------------------

Return to the hill. At every point on the hill, we can assign a height using a function 
$h(x, y)$. This function tells us the elevation at each position.

\medskip

Now suppose you stand at a point on the hill and choose a direction in which to walk. 
You might ask:

\begin{quote}
How steep is the hill in this direction?
\end{quote}

\medskip

This is exactly the kind of question a covector answers.

\medskip

The \textbf{gradient} of $h$,

\[
dh = \begin{bmatrix} \dfrac{\partial h}{\partial x} & \dfrac{\partial h}{\partial y} \end{bmatrix},
\]

is a covector. It takes a direction vector $\vec{v} = \begin{bmatrix} v_1 \\ v_2 \end{bmatrix}$ 
and returns a number:

\[
dh(\vec{v}) = \frac{\partial h}{\partial x}\,v_1 + \frac{\partial h}{\partial y}\,v_2 \in \mathbb{R}.
\]

\medskip

This number tells you how quickly the height changes if you move in the direction $\vec{v}$.

\medskip

\textbf{Key observation.} The vector $\vec{v}$ tells you \textit{which direction} you are moving 
on the hill. The covector $dh$ tells you \textit{how steeply} you are climbing in that direction.

\medskip

One describes motion; the other assigns a number to that motion.

\medskip

\section{Covectors in Coordinates}

In coordinates, a covector in $\mathbb{R}^n$ is represented as a \textbf{row vector}:

\[
f = \begin{bmatrix} f_1 & f_2 & \cdots & f_n \end{bmatrix},
\]

while an ordinary vector is represented as a \textbf{column vector}:

\[
\vec{v} =
\begin{bmatrix}
v_1 \\
v_2 \\
\vdots \\
v_n
\end{bmatrix}.
\]

\medskip

The action of $f$ on $\vec{v}$ is then simply matrix multiplication:

\[
f(\vec{v}) =
\begin{bmatrix} f_1 & f_2 & \cdots & f_n \end{bmatrix}
\begin{bmatrix} v_1 \\ v_2 \\ \vdots \\ v_n \end{bmatrix}
= f_1 v_1 + f_2 v_2 + \cdots + f_n v_n.
\]

\medskip

So a covector takes the components of a vector and combines them into a single number.

% ---------------------------------------------------------
\section{Vectors vs.\ Covectors}
% ---------------------------------------------------------

We can now state the fundamental distinction precisely:

\medskip

\begin{center}
\begin{tabular}{l@{\hspace{1.5cm}}l@{\hspace{1.5cm}}l}
\toprule
& \textbf{Vector} & \textbf{Covector} \\
\midrule
What it is   & Direction or displacement  & Linear function on vectors \\
Representation & Column vector        & Row vector \\
What it does  & Describes motion       & Assigns a number to motion \\
Output     & A vector           & A real number \\
\bottomrule
\end{tabular}
\end{center}

\medskip

\textbf{Key idea.} Vectors tell us how to move. Covectors tell us something about that movement by assigning a number to it.
% ---------------------------------------------------------
\section{How Covectors Transform}
% ---------------------------------------------------------

This is where the distinction becomes essential. When we change basis, vectors 
and covectors do \textit{not} transform in the same way --- and there is a precise 
reason why.

\medskip

\textbf{Key principle.} The number produced by a covector acting on a vector must 
not depend on how we describe the vector. It is a physical or geometric quantity, 
and so it must remain the same in every basis.

\medskip

Suppose we change basis using the matrix $P$. A vector $\vec{v}$ expressed in the 
new basis has components related to the old ones by:

\[
\vec{v}_{\,\text{new}} = P^{-1}\vec{v}_{\,\text{old}}.
\]

\medskip

Now consider the number $f(\vec{v})$. This is a real quantity --- for example, the 
rate of change of height on the hill in a given direction --- and it must be the same 
no matter which basis we use.

\medskip

This requirement forces the covector to transform in the opposite way.

\medskip

In the old basis:
\[
f(\vec{v}) = f_{\,\text{old}}\,\vec{v}_{\,\text{old}}.
\]

In the new basis, substituting $\vec{v}_{\,\text{old}} = P\,\vec{v}_{\,\text{new}}$:
\[
f(\vec{v}) = f_{\,\text{old}}\,P\,\vec{v}_{\,\text{new}}.
\]

\medskip

For this to have the same form in the new basis, we must define:
\[
f_{\,\text{new}} = f_{\,\text{old}}\,P.
\]

Rearranging:
\[
f_{\,\text{old}} = f_{\,\text{new}}\,P^{-1}.
\]
So that:
\[
f(\vec{v}) = f_{\,\text{new}}\,\vec{v}_{\,\text{new}}. \quad \checkmark
\]




\medskip

\textbf{Key idea.} Since the vector components change, the covector must adjust in 
exactly the opposite way so that the final number remains unchanged.

\medskip

The transformation rules are therefore:

\medskip

\begin{center}
\begin{tabular}{l@{\hspace{2cm}}c}
\toprule
& \textbf{Transformation rule} \\
\midrule
Vector components  & multiply by $P^{-1}$ \\
Covector components & multiply by $P$ \\
\bottomrule
\end{tabular}
\end{center}

\medskip

Because they transform by inverse matrices relative to one another, vectors and 
covectors are said to transform \textbf{contravariantly} and \textbf{covariantly} 
respectively. This is the origin of the prefix \textit{co} in covector.

% ---------------------------------------------------------
\section{Index Notation}
% ---------------------------------------------------------

In more advanced mathematics and physics, the distinction between vectors and 
covectors is built directly into the notation.

\medskip

Vector components are written with \textbf{upper indices}:
\[
\vec{v} = (v^1, v^2, \ldots, v^n),
\]

while covector components are written with \textbf{lower indices}:
\[
f = (f_1, f_2, \ldots, f_n).
\]

\medskip

This difference is not just stylistic --- it reflects how these objects transform 
when we change basis.

\medskip

The action of a covector on a vector is written as:
\[
f(\vec{v}) = f_i\, v^i.
\]

\medskip

This expression represents a sum over the index $i$:
\[
f_i\, v^i = f_1 v^1 + f_2 v^2 + \cdots + f_n v^n.
\]

\medskip

This is the \textbf{Einstein summation convention}: whenever an index appears 
once upstairs and once downstairs, it is automatically summed over.

\medskip

\textbf{Key idea.} The placement of indices is not cosmetic --- it encodes how 
each object transforms, and ensures that expressions like $f_i v^i$ produce a 
scalar that is independent of the choice of basis.

% ---------------------------------------------------------
\section{Why This Matters}
% ---------------------------------------------------------

The distinction between vectors and covectors is not merely technical --- it is 
essential.

\medskip

In physics, quantities such as velocities, displacements, and forces are vectors, 
while gradients and certain forms of momentum are covectors. Treating these as 
interchangeable can lead to expressions that depend on the coordinate system, and 
therefore do not represent physical reality.

\medskip

More fundamentally, covectors are the natural objects that act on vectors to 
produce scalars. This pairing,
\[
f_i v^i,
\]
gives a number that is independent of the choice of basis.

\medskip

At this point, a natural question arises:

\begin{quote}
Can we produce such a number using only vectors, without explicitly introducing a covector?
\end{quote}

\medskip

In other words, given two vectors, is there a natural way to assign a number that 
captures how they relate to one another?

\medskip

The answer is yes. This leads us to the concept of an \textbf{inner product}, which 
provides a way to measure lengths, angles, and alignment between vectors.

\medskip

In the next section, we introduce the inner product and explore how it allows 
vectors themselves to generate scalars, and how it connects vectors and covectors 
in a deeper way.

\chapter{The Inner Product}
% ---------------------------------------------------------

In the previous section, we saw that a covector acting on a vector produces a number:
\[
f_i v^i.
\]

This naturally leads to the question: can we assign a number to a pair of vectors 
directly, without introducing a separate covector?

\medskip

The answer is yes. This is done using an \textbf{inner product}.

\medskip

An inner product is a rule that takes two vectors and returns a real number:
\[
\langle \vec{u}, \vec{v} \rangle \in \mathbb{R}.
\]

\medskip

It satisfies the following properties for all vectors $\vec{u}, \vec{v}, \vec{w}$ 
and scalar $c$:
\begin{itemize}
  \item Linearity in each argument:
  \[
  \langle \vec{u} + \vec{w}, \vec{v} \rangle = \langle \vec{u}, \vec{v} \rangle + \langle \vec{w}, \vec{v} \rangle,
  \qquad
  \langle c\vec{u}, \vec{v} \rangle = c\,\langle \vec{u}, \vec{v} \rangle,
  \]
  \item Symmetry:
  \[
  \langle \vec{u}, \vec{v} \rangle = \langle \vec{v}, \vec{u} \rangle,
  \]
  \item Positivity:
  \[
  \langle \vec{v}, \vec{v} \rangle \geq 0,
  \quad \text{with equality only if } \vec{v} = \vec{0}.
  \]
\end{itemize}

\medskip

\textbf{Geometric meaning.} The inner product measures how two vectors relate to 
each other. It captures both their lengths and how aligned they are.

\medskip

In $\mathbb{R}^2$, the standard inner product is:
\[
\langle \vec{u}, \vec{v} \rangle = u^1 v^1 + u^2 v^2.
\]

\medskip

\section{Connection to the Hill}

Return to the hill. Suppose you choose a direction $\vec{v}$ in which to walk, and 
another direction $\vec{u}$.

\medskip

The inner product $\langle \vec{u}, \vec{v} \rangle$ tells you how much these two 
directions align with each other.

\begin{itemize}
  \item If the value is large and positive, the directions are closely aligned.
  \item If it is zero, the directions are perpendicular.
  \item If it is negative, the directions point in opposite directions.
\end{itemize}

\medskip

So while a vector tells you how to move, the inner product tells you how two 
movements compare.

\medskip

\section{From Vectors to Covectors}

An important consequence of the inner product is that it allows us to associate a 
covector to every vector.

\medskip

Given a vector $\vec{u}$, we can define a function:
\[
f_{\vec{u}}(\vec{v}) = \langle \vec{u}, \vec{v} \rangle.
\]

\medskip

This function takes a vector $\vec{v}$ and returns a number, and it is linear in 
$\vec{v}$. Therefore, it is a covector.

\medskip

\textbf{Key idea.} The inner product provides a bridge between vectors and 
covectors: it allows us to turn vectors into objects that measure other vectors.

\end{document}
